\section{Conclusions}

An auditable global FPV assessment must keep water-body identity, resource assumptions, screening evidence and project deployability separate. After restoring 1{,}245 lakes removed without duplicate evidence, the v117 reference inventory contains 199{,}976 large lakes and mapped reservoirs. Its reference L0 resource is 16{,}713~TWh~yr$^{-1}$, and the ice screen retains 10{,}995~TWh~yr$^{-1}$ (65.79\%).

Public conservation, population and 1991--2018 surface-area evidence support L2 intervals of 4{,}190--4{,}805~TWh~yr$^{-1}$ under the core definition and 2{,}914--3{,}404~TWh~yr$^{-1}$ under the conservative definition. Bounded road and mapped-power-line proximity reduce these to partial-L3 intervals of 1{,}595--3{,}508 and 1{,}101--2{,}494~TWh~yr$^{-1}$. Coverage choice controls the physical resource, ice evidence controls much of the lake contraction, and infrastructure mapping controls most of the partial-L3 range.

These endpoints have different decision roles. Lower values define higher-confidence global screening sets, whereas upper values identify where new observations could alter screening decisions. Complete suitability under the published 1991--2020 criterion and complete deployable potential remain unavailable. The next step is not another global discount factor, but targeted surface-area, grid-capacity, bathymetric, environmental and project-engineering evidence for the candidate portfolios identified here.
